\chapter{Compiler Architecture}
\label{ch:compiler}

This chapter describes the architecture and implementation of the \haira{} compiler, updated to reflect \haira{}'s identity as an agentic orchestration language.

\section{Overview}

The \haira{} compiler (\code{hairac}) transforms source code into native executables. It is written in Rust and uses Cranelift for code generation. Agentic constructs --- \keyword{provider}, \keyword{tool}, \keyword{agent}, and \keyword{workflow} with \code{@} decorator triggers --- are first-class elements of the compilation pipeline.

The compiler pipeline consists of the following stages:

\begin{lstlisting}[language={},numbers=none]
Source (.haira)
       |
       v
   +--------+
   | Lexer  |  --> Tokens
   +--------+
       |
       v
   +--------+
   | Parser |  --> AST
   +--------+
       |
       v
   +----------+
   | Resolver |  --> Resolved AST (imports, symbols, providers)
   +----------+
       |
       v
   +-------------+
   | Type Checker|  --> Typed AST (tools, agents, workflows validated)
   +-------------+
       |
       v
   +----------+
   | Lowering |  --> HIR (High-level IR)
   +----------+
       |
       v
   +----------+
   | Codegen  |  --> Cranelift IR --> Native binary
   +----------+
\end{lstlisting}

\begin{note}
There is no ``AI Phase'' in the pipeline. Previous versions of \haira{} included a compile-time AI code generation stage using the \keyword{ai} keyword and a Canonical Intermediate Representation (CIR). These have been removed. \haira{} is now a language for \emph{orchestrating} AI agents at runtime, not for generating code with AI at compile time.
\end{note}

\section{Lexer}

The lexer (scanner) converts source text into a stream of tokens.

\subsection{Token Types}

\begin{lstlisting}[language={}]
enum TokenKind {
    // Literals
    IntLit(i64),
    FloatLit(f64),
    StringLit(String),
    TripleQuoteStringLit(String),   // """..."""
    BoolLit(bool),

    // Identifiers and keywords
    Ident(String),
    Keyword(Keyword),

    // Operators
    Plus, Minus, Star, Slash, Percent,
    Eq, EqEq, NotEq, Lt, Gt, LtEq, GtEq,
    And, Or, Not, AndAnd, OrOr, Bang,
    PlusEq, MinusEq, StarEq, SlashEq,
    Pipe, DotDot, DotDotEq, Arrow, FatArrow,
    LeftArrow, Question,
    At,                              // @ (decorator prefix)

    // Delimiters
    LParen, RParen, LBracket, RBracket, LBrace, RBrace,
    Comma, Colon, Dot, Semicolon,

    // Special
    Newline, Eof, Error(String),
}
\end{lstlisting}

\subsection{Keywords}

The keyword enum includes the four agentic keywords:

\begin{lstlisting}[language={}]
enum Keyword {
    // General-purpose keywords
    And, As, Break, Catch, Chan, Continue, Defer,
    Else, Enum, False, Fn, For, From, If, Import,
    In, Match, Nil, Not, Or, Return, Select, Spawn,
    Struct, True, Try, Type, While,

    // Agentic keywords
    Agent,
    Provider,
    Tool,
    Workflow,
}
\end{lstlisting}

\begin{note}
The \keyword{ai} keyword from previous versions is no longer part of the language. It is not reserved and may be used as an identifier.
\end{note}

\subsection{Lexer Implementation}

The lexer is a hand-written scanner that operates on UTF-8 source text:

\begin{lstlisting}[language={}]
struct Lexer<'a> {
    source: &'a str,
    chars: Peekable<CharIndices<'a>>,
    line: usize,
    column: usize,
}

impl<'a> Lexer<'a> {
    fn next_token(&mut self) -> Token {
        self.skip_whitespace();

        match self.peek_char() {
            None => Token::new(TokenKind::Eof, self.span()),
            Some(c) => match c {
                '0'..='9' => self.scan_number(),
                '"' => self.scan_string_or_triple_quote(),
                'a'..='z' | 'A'..='Z' | '_' => self.scan_identifier(),
                '@' => self.single_token(TokenKind::At),
                '+' => self.scan_operator(TokenKind::Plus, TokenKind::PlusEq),
                // ... other cases
            }
        }
    }

    fn scan_string_or_triple_quote(&mut self) -> Token {
        if self.match_prefix("\"\"\"") {
            self.scan_triple_quote_string()
        } else {
            self.scan_string()
        }
    }
}
\end{lstlisting}

The \code{scan\_string\_or\_triple\_quote} method distinguishes between regular strings (\code{"..."}) and triple-quoted strings (\code{"""..."""}) by checking for three consecutive quote characters.

\section{Parser}

The parser constructs an Abstract Syntax Tree (AST) from the token stream using recursive descent with Pratt parsing for expressions.

\subsection{AST Nodes}

The core expression and statement nodes remain unchanged from the base language:

\begin{lstlisting}[language={}]
enum Expr {
    Literal(Literal),
    Ident(Ident),
    Binary(Box<Expr>, BinOp, Box<Expr>),
    Unary(UnaryOp, Box<Expr>),
    Call(Box<Expr>, Vec<Expr>),
    FieldAccess(Box<Expr>, Ident),
    Index(Box<Expr>, Box<Expr>),
    If(Box<Expr>, Block, Option<Block>),
    Match(Box<Expr>, Vec<MatchArm>),
    Closure(Vec<Param>, Option<Type>, Block),
    Array(Vec<Expr>),
    Map(Vec<(Expr, Expr)>),
    Struct(Ident, Vec<(Ident, Expr)>),
    Range(Box<Expr>, Box<Expr>, bool),
    Pipe(Box<Expr>, Box<Expr>),
}

enum Stmt {
    VarDecl(Vec<Ident>, Option<Type>, Vec<Expr>),
    Assignment(Vec<LValue>, AssignOp, Vec<Expr>),
    Expr(Expr),
    If(Expr, Block, Option<Block>),
    For(Option<Ident>, Ident, Expr, Block),
    While(Expr, Block),
    Match(Expr, Vec<MatchArm>),
    Return(Option<Vec<Expr>>),
    Break(Option<Ident>),
    Continue(Option<Ident>),
    Defer(Box<Stmt>),
    Spawn(SpawnTarget),
    Block(Block),
}
\end{lstlisting}

\subsection{Agentic AST Nodes}

Four new top-level declaration nodes support the agentic constructs:

\begin{lstlisting}[language={}]
struct ProviderDecl {
    name: Ident,
    fields: Vec<(Ident, Expr)>,
}

struct ToolDecl {
    name: Ident,
    params: Vec<Param>,
    return_type: Option<Type>,
    description: String,        // from triple-quoted string
    body: Block,
}

struct AgentDecl {
    name: Ident,
    fields: Vec<(Ident, AgentFieldValue)>,
}

struct WorkflowDecl {
    decorators: Vec<Decorator>,
    name: Ident,
    params: Vec<Param>,
    return_type: Option<Type>,
    body: Block,
}

struct Decorator {
    name: Ident,                // e.g., "webhook", "cron", "event"
    args: Vec<Expr>,            // e.g., "/summarize", "0 9 * * *"
}
\end{lstlisting}

\subsection{Parsing Agentic Constructs}

The recursive descent parser is extended to handle the new keywords at the top level:

\begin{lstlisting}[language={}]
impl Parser<'_> {
    fn parse_top_level(&mut self) -> Result<Decl, ParseError> {
        match self.peek_keyword() {
            Some(Keyword::Fn)       => self.parse_fn_decl(),
            Some(Keyword::Struct)   => self.parse_struct_decl(),
            Some(Keyword::Enum)     => self.parse_enum_decl(),
            Some(Keyword::Type)     => self.parse_type_decl(),
            Some(Keyword::Import)   => self.parse_import(),
            Some(Keyword::Provider) => self.parse_provider_decl(),
            Some(Keyword::Tool)     => self.parse_tool_decl(),
            Some(Keyword::Agent)    => self.parse_agent_decl(),
            Some(Keyword::Workflow) => self.parse_workflow_decl(vec![]),
            _ if self.check(TokenKind::At) => {
                let decorators = self.parse_decorators()?;
                self.expect_keyword(Keyword::Workflow)?;
                self.parse_workflow_decl(decorators)
            }
            _ => Err(ParseError::UnexpectedToken(self.peek().clone())),
        }
    }

    fn parse_decorators(&mut self) -> Result<Vec<Decorator>, ParseError> {
        let mut decorators = Vec::new();
        while self.check(TokenKind::At) {
            self.advance(); // consume @
            let name = self.parse_ident()?;
            let args = if self.check(TokenKind::LParen) {
                self.parse_call_args()?
            } else {
                vec![]
            };
            decorators.push(Decorator { name, args });
        }
        Ok(decorators)
    }

    fn parse_tool_decl(&mut self) -> Result<Decl, ParseError> {
        self.expect_keyword(Keyword::Tool)?;
        let name = self.parse_ident()?;
        let params = self.parse_params()?;
        let return_type = self.parse_optional_return_type()?;
        self.expect(TokenKind::LBrace)?;
        let description = self.expect_triple_quote_string()?;
        let body = self.parse_block_body()?;
        Ok(Decl::Tool(ToolDecl { name, params, return_type, description, body }))
    }
}
\end{lstlisting}

\begin{implementation}
The decorator parser greedily consumes all consecutive \code{@}-prefixed lines before the \keyword{workflow} keyword. This means decorators must appear immediately before the workflow declaration with no intervening blank lines or other declarations.
\end{implementation}

\section{Resolver}

The resolver performs name resolution, import handling, and symbol binding. It is extended to handle agentic constructs.

\subsection{Symbol Table}

The symbol table is extended with new symbol kinds for agentic constructs:

\begin{lstlisting}[language={}]
struct SymbolTable {
    scopes: Vec<Scope>,
    current: usize,
}

struct Scope {
    symbols: HashMap<String, Symbol>,
    parent: Option<usize>,
}

enum Symbol {
    Variable { ty: TypeId, mutable: bool },
    Function { sig: FunctionSig },
    Type { def: TypeDef },
    Module { exports: HashMap<String, Symbol> },
    Provider { fields: ProviderFields },
    Tool { sig: ToolSig, description: String },
    Agent { fields: AgentFields },
    Workflow { sig: WorkflowSig, triggers: Vec<Trigger> },
}
\end{lstlisting}

\subsection{Provider Resolution}

The resolver validates provider declarations:

\begin{itemize}
    \item The \code{model} field is required.
    \item The \code{api\_key} field, if present, must be a string expression (typically an \code{env()} call).
    \item The \code{backend} field, if present, must be a recognized backend identifier (\code{"ollama"}, \code{"llama.cpp"}).
    \item Provider names are registered in the module scope for reference by agents.
\end{itemize}

\begin{lstlisting}[language={}]
fn resolve_provider(&mut self, decl: &ProviderDecl) -> Result<(), ResolveError> {
    // Ensure required fields are present
    if !decl.has_field("model") {
        return Err(ResolveError::MissingField {
            construct: "provider",
            name: decl.name.clone(),
            field: "model",
        });
    }

    // Register in symbol table
    self.bind_symbol(
        &decl.name,
        Symbol::Provider { fields: self.extract_provider_fields(decl)? },
    )
}
\end{lstlisting}

\subsection{Tool Registration}

The resolver processes tool declarations and prepares metadata for schema generation:

\begin{itemize}
    \item The tool name is registered in the module scope.
    \item Parameter types are recorded for JSON schema generation.
    \item The triple-quoted description is stored as metadata.
    \item Default parameter values are validated.
\end{itemize}

\begin{lstlisting}[language={}]
fn resolve_tool(&mut self, decl: &ToolDecl) -> Result<(), ResolveError> {
    // Description is mandatory (enforced by parser, double-checked here)
    if decl.description.is_empty() {
        return Err(ResolveError::MissingToolDescription(decl.name.clone()));
    }

    // Resolve parameter types and body
    self.push_scope();
    for param in &decl.params {
        self.resolve_type(&param.ty)?;
        self.bind_local(&param.name, &param.ty)?;
    }
    self.resolve_block(&decl.body)?;
    self.pop_scope();

    // Register tool symbol
    self.bind_symbol(
        &decl.name,
        Symbol::Tool {
            sig: self.build_tool_sig(decl)?,
            description: decl.description.clone(),
        },
    )
}
\end{lstlisting}

\subsection{Agent Validation}

The resolver verifies that agent declarations reference valid providers and tools:

\begin{itemize}
    \item The \code{model} field must reference a declared provider.
    \item Each entry in the \code{tools} list must reference a declared tool.
    \item The \code{system} field must be a string expression.
    \item Memory configuration, if present, must use a recognized memory type.
\end{itemize}

\begin{lstlisting}[language={}]
fn resolve_agent(&mut self, decl: &AgentDecl) -> Result<(), ResolveError> {
    // Validate model references a declared provider
    let model_field = decl.get_field("model")
        .ok_or(ResolveError::MissingField {
            construct: "agent", name: decl.name.clone(), field: "model",
        })?;
    self.expect_provider_ref(model_field)?;

    // Validate tool references
    if let Some(tools) = decl.get_field("tools") {
        for tool_name in tools.as_ident_list()? {
            self.expect_tool_ref(&tool_name)?;
        }
    }

    // Register agent symbol
    self.bind_symbol(
        &decl.name,
        Symbol::Agent { fields: self.extract_agent_fields(decl)? },
    )
}
\end{lstlisting}

\subsection{Workflow Trigger Validation}

The resolver validates workflow decorators:

\begin{itemize}
    \item \code{@webhook} requires a string literal path argument.
    \item \code{@cron} requires a valid cron expression string.
    \item \code{@event} requires a string literal event name.
    \item \code{@websocket} requires a string literal path argument.
    \item \code{@manual} takes no arguments.
    \item Unrecognized decorator names produce an error.
\end{itemize}

\begin{lstlisting}[language={}]
fn resolve_workflow(&mut self, decl: &WorkflowDecl) -> Result<(), ResolveError> {
    // Validate each decorator
    let mut triggers = Vec::new();
    for decorator in &decl.decorators {
        let trigger = match decorator.name.as_str() {
            "webhook"   => self.validate_webhook_trigger(decorator)?,
            "websocket" => self.validate_websocket_trigger(decorator)?,
            "cron"      => self.validate_cron_trigger(decorator)?,
            "event"     => self.validate_event_trigger(decorator)?,
            "manual"    => self.validate_manual_trigger(decorator)?,
            unknown     => return Err(ResolveError::UnknownTrigger {
                name: unknown.to_string(),
                span: decorator.span,
            }),
        };
        triggers.push(trigger);
    }

    // Resolve params and body
    self.push_scope();
    for param in &decl.params {
        self.resolve_type(&param.ty)?;
        self.bind_local(&param.name, &param.ty)?;
    }
    self.resolve_block(&decl.body)?;
    self.pop_scope();

    // Register workflow symbol
    self.bind_symbol(
        &decl.name,
        Symbol::Workflow {
            sig: self.build_workflow_sig(decl)?,
            triggers,
        },
    )
}
\end{lstlisting}

\section{Type Checker}

The type checker verifies type correctness and performs type inference. It is extended to handle the types and patterns introduced by agentic constructs.

\subsection{Type Representation}

The type system includes a new \type{stream<T>} type for streaming agent responses:

\begin{lstlisting}[language={}]
enum Type {
    // Primitives
    Int, I8, I16, I32, I64,
    UInt, U8, U16, U32, U64,
    Float, F32, F64,
    Bool,
    String,

    // Composites
    Array(Box<Type>),
    Map(Box<Type>, Box<Type>),
    Tuple(Vec<Type>),
    Option(Box<Type>),

    // Named
    Struct(StructId),
    Enum(EnumId),

    // Functions
    Function(Vec<Type>, Box<Type>),

    // Generics
    TypeVar(TypeVarId),
    Generic(GenericId, Vec<Type>),

    // Streaming
    Stream(Box<Type>),          // stream<T>

    // Special
    Never,
    Error,
}
\end{lstlisting}

\subsection{Agent Method Type Checking}

Agent methods have specific type signatures that the type checker enforces:

\begin{itemize}
    \item \code{Agent.ask(prompt: string, session: string?) -> (string, Error?)} --- Text in, text out.
    \item \code{Agent.run(input: T, session: string?) -> (U, Error?)} --- Structured in, structured out. The output type \code{U} is inferred from the left-side type annotation at the call site.
    \item \code{Agent.stream(prompt: string, session: string?) -> stream<string>} --- Text in, streaming out.
\end{itemize}

\begin{lstlisting}[language={}]
fn check_agent_method_call(
    &mut self,
    agent: &Symbol,
    method: &str,
    args: &[Expr],
    expected_return: Option<&Type>,
) -> Result<Type, TypeError> {
    match method {
        "ask" => {
            self.check_arg_type(&args[0], &Type::String)?;
            Ok(Type::Tuple(vec![Type::String, Type::Option(Box::new(Type::Error))]))
        }

        "run" => {
            let input_ty = self.infer_expr(&args[0])?;
            self.check_has_json_schema(&input_ty)?;

            // Output type inferred from left-side annotation
            let output_ty = expected_return
                .ok_or(TypeError::MissingTypeAnnotation {
                    hint: "agent.run() requires a type annotation on the left side",
                })?;
            self.check_has_json_schema(output_ty)?;

            Ok(Type::Tuple(vec![
                output_ty.clone(),
                Type::Option(Box::new(Type::Error)),
            ]))
        }

        "stream" => {
            self.check_arg_type(&args[0], &Type::String)?;
            Ok(Type::Stream(Box::new(Type::String)))
        }

        _ => Err(TypeError::NoSuchMethod {
            ty: "agent",
            method: method.to_string(),
        }),
    }
}
\end{lstlisting}

\subsection{Left-Side Type Annotation Inference}

For \code{Agent.run()}, the output type is determined by the type annotation on the left side of the assignment:

\begin{lstlisting}
// The compiler infers that Researcher.run() should produce ResearchResult
research: ResearchResult, err = Researcher.run(ResearchRequest{
    topic: "quantum computing"
})
\end{lstlisting}

The type checker extracts \type{ResearchResult} from the left-side annotation and uses it as the expected output type. It then validates that both \type{ResearchRequest} and \type{ResearchResult} are struct types with JSON-serializable fields, so that the runtime can generate and validate JSON schemas.

\subsection{Workflow Input/Output Types}

Workflow parameters define the input type. For triggered workflows, the compiler validates that parameter and return types are JSON-serializable:

\begin{lstlisting}[language={}]
fn check_workflow(&mut self, decl: &WorkflowDecl) -> Result<(), TypeError> {
    // For triggered workflows, params must be JSON-serializable
    if !decl.decorators.is_empty() {
        for param in &decl.params {
            self.check_json_serializable(&param.ty)?;
        }
        if let Some(ret_ty) = &decl.return_type {
            self.check_json_serializable(ret_ty)?;
        }
    }

    // Check body
    self.push_scope();
    for param in &decl.params {
        self.bind_local(&param.name, &param.ty);
    }
    let body_ty = self.check_block(&decl.body)?;

    // Verify return type matches
    if let Some(ret_ty) = &decl.return_type {
        self.unify(&body_ty, ret_ty)?;
    }
    self.pop_scope();

    Ok(())
}
\end{lstlisting}

\subsection{Tool Auto-Schema Generation}

The type checker produces JSON schema metadata for each tool declaration. This metadata is embedded in the compiled binary and used by the runtime when registering tools with LLM providers:

\begin{lstlisting}[language={}]
fn generate_tool_schema(&self, decl: &ToolDecl) -> ToolSchema {
    ToolSchema {
        name: decl.name.clone(),
        description: decl.description.clone(),
        parameters: JsonSchema::Object {
            properties: decl.params.iter().map(|p| {
                (p.name.clone(), self.type_to_json_schema(&p.ty))
            }).collect(),
            required: decl.params.iter()
                .filter(|p| p.default.is_none())
                .map(|p| p.name.clone())
                .collect(),
        },
    }
}
\end{lstlisting}

The type-to-schema mapping is:

\begin{table}[h]
\centering
\begin{tabular}{ll}
\toprule
\textbf{\haira{} Type} & \textbf{JSON Schema Type} \\
\midrule
\type{int} & \code{"integer"} \\
\type{float} & \code{"number"} \\
\type{bool} & \code{"boolean"} \\
\type{string} & \code{"string"} \\
\type{[T]} & \code{"array"} with \code{items} \\
\type{Struct} & \code{"object"} with \code{properties} \\
\type{T?} & adds \code{"null"} to type union \\
\bottomrule
\end{tabular}
\caption{Type-to-JSON-schema mapping}
\end{table}

\section{Lowering}

Lowering transforms the typed AST into HIR (High-level IR). Agentic constructs are lowered to calls into the \haira{} runtime library.

\subsection{HIR Structure}

\begin{lstlisting}[language={}]
struct HirModule {
    functions: Vec<HirFunction>,
    types: Vec<HirType>,
    constants: Vec<HirConst>,
    providers: Vec<HirProvider>,
    tools: Vec<HirTool>,
    agents: Vec<HirAgent>,
    workflows: Vec<HirWorkflow>,
}

struct HirFunction {
    name: Symbol,
    params: Vec<HirParam>,
    return_type: HirType,
    body: Vec<HirStmt>,
    locals: Vec<HirLocal>,
}

enum HirStmt {
    Assign(LocalId, HirExpr),
    Store(HirPlace, HirExpr),
    Call(Option<LocalId>, FnId, Vec<HirExpr>),
    RuntimeCall(Option<LocalId>, RuntimeFn, Vec<HirExpr>),
    If(HirExpr, BlockId, Option<BlockId>),
    Loop(BlockId),
    Break(Option<LoopId>),
    Continue(Option<LoopId>),
    Return(Option<HirExpr>),
}
\end{lstlisting}

\subsection{Provider Lowering}

A \keyword{provider} declaration is lowered to a runtime configuration call:

\begin{lstlisting}[language={}]
// Source:
//   provider anthropic {
//       api_key: env("ANTHROPIC_API_KEY")
//       model: "claude-sonnet-4-20250514"
//   }

// Lowered to:
HirStmt::RuntimeCall(
    Some(provider_local),
    RuntimeFn::RegisterProvider,
    vec![
        HirExpr::StringLit("anthropic"),
        HirExpr::Map(vec![
            ("api_key", HirExpr::Call(env_fn, vec![HirExpr::StringLit("ANTHROPIC_API_KEY")])),
            ("model", HirExpr::StringLit("claude-sonnet-4-20250514")),
        ]),
    ],
)
\end{lstlisting}

\subsection{Tool Lowering}

A \keyword{tool} declaration is lowered to a runtime tool registration that includes the compiled function body and the auto-generated JSON schema:

\begin{lstlisting}[language={}]
// Source:
//   tool search(query: string, max_results: int = 5) -> [SearchResult] {
//       """Search the web for information"""
//       // ... body ...
//   }

// Lowered to:
// 1. The tool body becomes a regular HirFunction
HirFunction {
    name: Symbol("__tool_search"),
    params: vec![param_query, param_max_results],
    return_type: array_of_search_result,
    body: /* lowered body */,
}

// 2. A runtime registration call with the JSON schema
HirStmt::RuntimeCall(
    None,
    RuntimeFn::RegisterTool,
    vec![
        HirExpr::StringLit("search"),
        HirExpr::FnRef(Symbol("__tool_search")),
        HirExpr::StaticData(tool_json_schema),  // compiled schema bytes
    ],
)
\end{lstlisting}

\subsection{Agent Lowering}

An \keyword{agent} declaration is lowered to a runtime agent instantiation:

\begin{lstlisting}[language={}]
// Source:
//   agent Researcher {
//       model: anthropic
//       system: "You are a thorough researcher."
//       tools: [search, read_url]
//       temperature: 0.2
//       max_steps: 15
//   }

// Lowered to:
HirStmt::RuntimeCall(
    Some(agent_local),
    RuntimeFn::CreateAgent,
    vec![
        HirExpr::StringLit("Researcher"),
        HirExpr::ProviderRef("anthropic"),
        HirExpr::StringLit("You are a thorough researcher."),
        HirExpr::Array(vec![
            HirExpr::ToolRef("search"),
            HirExpr::ToolRef("read_url"),
        ]),
        HirExpr::Map(vec![
            ("temperature", HirExpr::FloatLit(0.2)),
            ("max_steps", HirExpr::IntLit(15)),
        ]),
    ],
)
\end{lstlisting}

\subsection{Agent Method Call Lowering}

Agent method calls are lowered to runtime API calls:

\begin{lstlisting}[language={}]
// Agent.ask() --> runtime ask call
// answer, err = Researcher.ask("query")
HirStmt::RuntimeCall(
    Some(result_local),
    RuntimeFn::AgentAsk,
    vec![
        HirExpr::AgentRef("Researcher"),
        HirExpr::StringLit("query"),
        HirExpr::Nil,  // no session
    ],
)

// Agent.run() --> runtime run call with output schema
// result: T, err = Agent.run(input)
HirStmt::RuntimeCall(
    Some(result_local),
    RuntimeFn::AgentRun,
    vec![
        HirExpr::AgentRef("Agent"),
        input_expr,
        HirExpr::StaticData(input_schema),
        HirExpr::StaticData(output_schema),
        HirExpr::Nil,  // no session
    ],
)

// Agent.stream() --> runtime stream call
// for chunk in Agent.stream("prompt") { ... }
HirStmt::RuntimeCall(
    Some(stream_local),
    RuntimeFn::AgentStream,
    vec![
        HirExpr::AgentRef("Agent"),
        HirExpr::StringLit("prompt"),
        HirExpr::Nil,  // no session
    ],
)
\end{lstlisting}

\subsection{Workflow Lowering}

Workflows are lowered differently based on their triggers:

\begin{lstlisting}[language={}]
// @webhook("/summarize")
// workflow Summarizer(url: string) -> Summary { ... }

// Lowered to:
// 1. The workflow body becomes a regular function
HirFunction {
    name: Symbol("__workflow_Summarizer"),
    params: vec![param_url],
    return_type: summary_type,
    body: /* lowered body */,
}

// 2. A runtime trigger registration
HirStmt::RuntimeCall(
    None,
    RuntimeFn::RegisterWebhook,
    vec![
        HirExpr::StringLit("/summarize"),
        HirExpr::StringLit("POST"),
        HirExpr::FnRef(Symbol("__workflow_Summarizer")),
        HirExpr::StaticData(input_schema),
        HirExpr::StaticData(output_schema),
    ],
)
\end{lstlisting}

\begin{table}[h]
\centering
\begin{tabular}{ll}
\toprule
\textbf{Trigger} & \textbf{Runtime Registration} \\
\midrule
\code{@webhook} & \code{RuntimeFn::RegisterWebhook} --- HTTP handler \\
\code{@websocket} & \code{RuntimeFn::RegisterWebsocket} --- WebSocket handler \\
\code{@cron} & \code{RuntimeFn::RegisterCron} --- Cron job \\
\code{@event} & \code{RuntimeFn::RegisterEvent} --- Event listener \\
\code{@manual} & \code{RuntimeFn::RegisterManual} --- CLI-only \\
(none) & No registration --- callable as sub-workflow \\
\bottomrule
\end{tabular}
\caption{Workflow trigger lowering}
\end{table}

\subsection{Spawn Block Lowering}

A \keyword{spawn} block inside a workflow is lowered to parallel execution runtime calls:

\begin{lstlisting}[language={}]
// Source:
//   reviews = spawn {
//       Reviewer.run(req1)
//       Reviewer.run(req2)
//       Reviewer.run(req3)
//   }

// Lowered to:
HirStmt::RuntimeCall(
    Some(reviews_local),
    RuntimeFn::SpawnAll,
    vec![
        HirExpr::Array(vec![
            HirExpr::Closure(/* Reviewer.run(req1) */),
            HirExpr::Closure(/* Reviewer.run(req2) */),
            HirExpr::Closure(/* Reviewer.run(req3) */),
        ]),
    ],
)
\end{lstlisting}

\section{Code Generation}

Code generation uses Cranelift to produce native machine code.

\subsection{Cranelift Integration}

The code generation stage translates HIR into Cranelift IR, which is then compiled to native code. The process is the same as for non-agentic code --- function bodies, control flow, and expressions are translated directly:

\begin{lstlisting}[language={}]
fn codegen_function(&mut self, func: &HirFunction) -> Result<(), CodegenError> {
    let mut builder = FunctionBuilder::new(&mut self.ctx.func, &mut self.builder_ctx);

    // Create entry block
    let entry = builder.create_block();
    builder.append_block_params_for_function_params(entry);
    builder.switch_to_block(entry);

    // Generate parameters
    for (i, param) in func.params.iter().enumerate() {
        let val = builder.block_params(entry)[i];
        self.locals.insert(param.id, val);
    }

    // Generate body
    for stmt in &func.body {
        self.codegen_stmt(&mut builder, stmt)?;
    }

    builder.seal_all_blocks();
    builder.finalize();

    Ok(())
}
\end{lstlisting}

\subsection{Runtime Calls}

Agentic constructs compile to calls into the \haira{} runtime library (\code{libhaira-runtime}). The code generator emits standard function calls to runtime symbols:

\begin{lstlisting}[language={}]
fn codegen_runtime_call(
    &mut self,
    builder: &mut FunctionBuilder,
    func: RuntimeFn,
    args: &[HirExpr],
) -> Value {
    let runtime_fn = self.import_runtime_function(func);
    let arg_values: Vec<Value> = args.iter()
        .map(|a| self.codegen_expr(builder, a))
        .collect();
    let call = builder.ins().call(runtime_fn, &arg_values);
    builder.inst_results(call)[0]
}
\end{lstlisting}

The compiled binary links against \code{libhaira-runtime}, which provides the implementations for all agentic operations (LLM API calls, tool execution, workflow orchestration, etc.). See Section~\ref{sec:runtime} for details.

\begin{note}
All agentic behavior happens at \emph{runtime}, not compile time. The compiler's role is to validate, type-check, and generate schemas for agentic constructs, then emit calls into the runtime library. The compiled binary is standalone but requires network access to reach LLM providers.
\end{note}

\section{Runtime Library}
\label{sec:runtime}

The \haira{} runtime library (\code{libhaira-runtime}) provides the execution environment for agentic constructs. It is linked into every compiled \haira{} binary that uses providers, tools, agents, or workflows.

\subsection{Provider Management}

The runtime manages provider connections:

\begin{itemize}
    \item \textbf{API key resolution} --- Reads API keys from environment variables at startup.
    \item \textbf{Model selection} --- Routes requests to the correct LLM backend based on provider configuration.
    \item \textbf{Connection pooling} --- Maintains HTTP connection pools to provider APIs.
    \item \textbf{Rate limiting} --- Handles provider rate limits with automatic retry and backoff.
    \item \textbf{Local backends} --- Manages connections to local inference servers (Ollama, llama.cpp).
\end{itemize}

\subsection{Agent Execution}

The runtime implements the agent execution loop:

\begin{enumerate}
    \item Assemble the message payload (system prompt, conversation history, user message).
    \item Send the request to the LLM provider.
    \item If the LLM response contains tool calls, execute each tool and collect results.
    \item Send tool results back to the LLM.
    \item Repeat steps 2--4 until the LLM produces a final response or \code{max\_steps} is reached.
    \item For \code{.run()} calls, parse the LLM response as structured JSON and deserialize into the expected output type.
    \item For \code{.stream()} calls, yield tokens as they arrive from the LLM.
\end{enumerate}

\subsection{Workflow Orchestration}

The runtime provides infrastructure for workflow triggers:

\begin{itemize}
    \item \textbf{HTTP server} --- An embedded HTTP server (built on a lightweight async runtime) handles \code{@webhook} and \code{@websocket} triggers. Request bodies are deserialized into workflow parameter types; return values are serialized as JSON responses.
    \item \textbf{Cron scheduler} --- A built-in cron scheduler executes \code{@cron}-triggered workflows on schedule.
    \item \textbf{Event bus} --- An in-process event bus dispatches \code{@event}-triggered workflows when events are emitted via \code{haira.emit()}.
    \item \textbf{Manual runner} --- The CLI can invoke \code{@manual} workflows directly with JSON input.
\end{itemize}

\subsection{Session Management}

The runtime manages conversation memory scoped by session ID:

\begin{itemize}
    \item Each session maintains an isolated conversation history.
    \item \code{conversation(max\_turns: N)} memory trims history to the last N message pairs.
    \item \code{summary(max\_tokens: N)} memory periodically summarizes older messages to stay within the token budget.
    \item Sessions are stored in-memory by default, with optional persistence backends.
\end{itemize}

\subsection{Stream Handling}

The runtime provides async token streaming for \code{Agent.stream()}:

\begin{itemize}
    \item Server-sent events (SSE) from LLM providers are parsed incrementally.
    \item Tokens are yielded to the caller via an async iterator interface.
    \item For \code{@websocket} workflows returning \code{stream<string>}, tokens are forwarded directly to the WebSocket connection.
\end{itemize}

\subsection{JSON Schema Generation}

The runtime uses type metadata embedded by the compiler to:

\begin{itemize}
    \item Register tool schemas with LLM providers during tool-calling.
    \item Instruct agents to produce structured output matching a given schema (for \code{.run()} calls).
    \item Validate and deserialize incoming webhook request bodies.
    \item Serialize workflow return values as JSON responses.
\end{itemize}

\section{Project Structure}

The compiler and runtime are organized as a Rust workspace with the following crates:

\begin{lstlisting}[language={},numbers=none]
haira/
  Cargo.toml                  # Workspace root

  crates/
    haira-lexer/              # Lexical analysis
    haira-parser/             # Parsing (+ provider, tool, agent, workflow)
    haira-ast/                # AST definitions (+ agentic nodes)
    haira-resolver/           # Name resolution (+ provider/agent/tool resolution)
    haira-types/              # Type system (+ stream<T>)
    haira-typeck/             # Type checking
    haira-hir/                # High-level IR
    haira-codegen/            # Code generation (Cranelift)
    haira-runtime/            # Runtime library (agents, workflows, providers)
    haira-driver/             # Compilation driver
    haira-cli/                # CLI (build, run, serve, check)
    haira-lsp/                # Language server
\end{lstlisting}

\begin{note}
The \code{haira-ai}, \code{haira-cir}, and \code{haira-local-ai} crates from previous versions have been removed. Their responsibilities are replaced by \code{haira-runtime}, which handles all agentic execution at runtime rather than compile time.
\end{note}

Each crate has a focused responsibility:

\begin{table}[h]
\centering
\begin{tabular}{ll}
\toprule
\textbf{Crate} & \textbf{Responsibility} \\
\midrule
\code{haira-lexer} & Tokenizes source text \\
\code{haira-parser} & Builds AST from tokens \\
\code{haira-ast} & Defines all AST node types \\
\code{haira-resolver} & Resolves imports, symbols, providers, tools, agents \\
\code{haira-types} & Type definitions and type utilities \\
\code{haira-typeck} & Type checking and inference \\
\code{haira-hir} & High-level IR definitions \\
\code{haira-codegen} & Cranelift code generation \\
\code{haira-runtime} & Runtime library for agentic execution \\
\code{haira-driver} & Orchestrates the compilation pipeline \\
\code{haira-cli} & Command-line interface \\
\code{haira-lsp} & Language server protocol implementation \\
\bottomrule
\end{tabular}
\caption{Crate responsibilities}
\end{table}

\section{Build Process}

The \code{haira} CLI provides the following commands:

\begin{lstlisting}[language=bash,numbers=none]
# Compile to native binary
haira build main.haira

# Compile and run immediately
haira run main.haira

# Compile and start server (for trigger-based workflows)
haira serve main.haira

# Type-check without code generation
haira check main.haira

# Interactive REPL
haira repl
\end{lstlisting}

\subsection{Build Modes}

\begin{lstlisting}[language=bash,numbers=none]
# Debug build (default) --- fast compilation, includes debug info
haira build main.haira

# Release build --- optimized, slower compilation
haira build main.haira --release
\end{lstlisting}

\subsection{Server Mode}

The \code{haira serve} command compiles the program and starts the runtime server, which listens for webhook requests, WebSocket connections, cron schedules, and events:

\begin{lstlisting}[language=bash,numbers=none]
# Start server on default port (8080)
haira serve main.haira

# Start server on a specific port
haira serve main.haira --port 3000

# Start with environment variables
ANTHROPIC_API_KEY=sk-... haira serve main.haira
\end{lstlisting}

\subsection{Check Mode}

The \code{haira check} command runs the full pipeline up to and including type checking, but skips code generation. This is useful for fast feedback during development:

\begin{lstlisting}[language=bash,numbers=none]
haira check main.haira
\end{lstlisting}

It validates all agentic constructs --- provider fields, tool descriptions, agent references, workflow triggers, and type compatibility --- without producing a binary.

\section{Error Messages}

The compiler produces detailed, contextual error messages. Errors related to agentic constructs follow the same format as core language errors.

\subsection{Missing Tool Description}

\begin{lstlisting}[language={},numbers=none]
error[E0101]: tool is missing a description
  --> src/tools.haira:5:1
   |
5  | tool search(query: string) -> [Result] {
   | ^^^^^^^^^^^^^^^^^^^^^^^^^^^^^^^^^^^^^^^^ tool declaration
6  |     resp = http.get("https://api.search.com?q={query}")
   |     ---- expected a """...""" description as the first element
   |
   = help: add a triple-quoted description: """Search the web for information"""
   = note: tool descriptions are sent to the LLM and cannot be omitted
\end{lstlisting}

\subsection{Agent Referencing Undefined Provider}

\begin{lstlisting}[language={},numbers=none]
error[E0102]: undefined provider `openai`
  --> src/agents.haira:8:12
   |
8  |     model: openai
   |            ^^^^^^ not found in this scope
   |
   = help: did you mean `anthropic`?
   = note: declare a provider with: provider openai { model: "gpt-4o" ... }
\end{lstlisting}

\subsection{Workflow with Invalid Trigger}

\begin{lstlisting}[language={},numbers=none]
error[E0103]: unknown trigger `@daily`
  --> src/workflows.haira:12:1
   |
12 | @daily
   | ^^^^^^ not a recognized trigger
   |
   = help: valid triggers are: @webhook, @websocket, @cron, @event, @manual
   = help: for daily execution, use: @cron("0 0 * * *")
\end{lstlisting}

\subsection{Type Mismatch in Agent.run() Output}

\begin{lstlisting}[language={},numbers=none]
error[E0104]: type mismatch in agent.run() output
  --> src/main.haira:20:5
   |
20 |     result: string, err = Analyst.run(request)
   |             ^^^^^^ expected a struct type, found `string`
   |
   = note: agent.run() requires a struct type annotation for structured output
   = help: define a struct for the output:
           struct AnalysisResult { summary: string, score: float }
           result: AnalysisResult, err = Analyst.run(request)
\end{lstlisting}

\subsection{Agent Referencing Undefined Tool}

\begin{lstlisting}[language={},numbers=none]
error[E0105]: undefined tool `search_web`
  --> src/agents.haira:10:18
   |
10 |     tools: [search_web, read_url]
   |             ^^^^^^^^^^ not found in this scope
   |
   = help: did you mean `search`?
   = note: tools must be declared with the `tool` keyword before being
           referenced in an agent
\end{lstlisting}

\subsection{Invalid Cron Expression}

\begin{lstlisting}[language={},numbers=none]
error[E0106]: invalid cron expression
  --> src/workflows.haira:3:7
   |
3  | @cron("every day")
   |        ^^^^^^^^^ not a valid cron expression
   |
   = help: cron expressions use the format: minute hour day-of-month month day-of-week
   = help: example: @cron("0 9 * * *") for every day at 9 AM
\end{lstlisting}
